\subsection{Example: Counting words}

\begin{frame}[t,fragile]{Counting words: problem and helper functions}
    \begin{itemize}
        \item Goal: count words in a text.
        \item A new word starts when a letter is preceded by a non-letter.
        \item The first character must be handled separately.
    \end{itemize}

\begin{block}{Helper functions}
\begin{lstlisting}
  bool is_letter(char c) {
    return std::isalpha(static_cast<unsigned char>(c));
  }

  size_t is_new_word(char c1, char c2) {
    return !is_letter(c1) && is_letter(c2);
  }
\end{lstlisting}
\end{block}
\end{frame}

\begin{frame}[t,fragile]{Using an algorithm: the map-reduce approach}
    \begin{itemize}
        \item Map: check for new words in pairs of characters.
          \begin{itemize}
              \item Sequences: $(t_0, \ldots, t_{n-2})$ and $(t_1, \ldots, t_{n-1})$.
              \item Mapping: \cppcode{is_new_word(c1,c2)} applied to pairs of characters.
          \end{itemize}
        \item Reduce: \cppcode{std::plus} for adding results of the mapping.
    \end{itemize}

\begin{block}{Map-reduce approach}
\begin{lstlisting}
  size_t count_classic_stl(std::string_view text) {
    if (text.empty()) { return 0; }
    auto result =
        std::transform_reduce(std::begin(text), std::prev(std::end(text)), 
        std::next(std::begin(text)), 
        0ZU,
        std::plus{}, 
        is_new_word);
    return result + is_letter(text[0]);
  }
\end{lstlisting}
\end{block}
\end{frame}

\begin{frame}[t,fragile]{Using a parallel map-reduce with C++17}
    \begin{itemize}
        \item The same map-reduce approach, but with a parallel execution policy.
    \end{itemize}

\begin{block}{Parallel map-reduce approach}
\begin{lstlisting}
  size_t count_parstl(std::string_view text) {
    if (text.empty()) { return 0; }
    auto result = std::transform_reduce(std::execution::par, 
        std::begin(text), std::prev(std::end(text)),
        std::next(std::begin(text)), 
        0ZU, 
        std::plus{},
        is_new_word);
    return result + is_letter(text[0]);
  }
\end{lstlisting}
\end{block}
\end{frame}

\begin{frame}[t,fragile]{Using ranges and views}
\begin{itemize}
  \item Steps:
    \begin{itemize}
      \item Generate a view with the $n-1$ first characters.
      \item Generate a view with the $n-1$ last characters.
      \item Zip the two views to get a view of pairs of characters.
      \item Count the number of pairs that start a new word.
    \end{itemize}
\end{itemize}

\begin{block}{Using ranges and views}
\begin{lstlisting}
  size_t count_ranges(std::string_view text) {
    if (text.empty()) { return 0; }
    auto compute = std::ranges::views::zip(
      text | std::ranges::views::take(std::ranges::size(text) - 1),
      text | std::ranges::views::drop(1));
    auto result = std::ranges::count_if(compute, [](auto char_pair) {
      auto [c1, c2] = char_pair;
      return is_new_word(c1, c2);
    });
    return result + is_letter(text[0]);
  }
\end{lstlisting}
\end{block}
\end{frame}

\begin{frame}[t,fragile]{Using parallel ranges and views with C++26}
\begin{itemize}
  \item The same approach as before, but with a parallel execution policy.
  \item \textmark{Note:} The parallel execution policy is not yet supported by all compilers.
    \begin{itemize}
      \item Experimental preview in OneAPI DPL.
    \end{itemize}
\end{itemize}

\begin{block}{Using parallel ranges and views}
\begin{lstlisting}
  size_t count_pranges(std::string_view text) {
    namespace stdx = oneapi::dpl;

    if (text.empty()) { return 0; }
    auto compute =
        std::ranges::views::zip(text | std::ranges::views::take(std::ranges::size(text) - 1),
                                text | std::ranges::views::drop(1));
    auto result = stdx::ranges::count_if(stdx::execution::par_unseq, compute, [](auto char_pair) {
      auto [c1, c2] = char_pair;
      return is_new_word(c1, c2);
    });
    return result + is_letter(text[0]);
  }
\end{lstlisting}
\end{block}
\end{frame}

\begin{frame}[t,fragile]{Performance results}
\begin{itemize}
  \item Counting the words in \textgood{``Don Quixote''} by Miguel de Cervantes.
  \item \textmark{Number of characters}: 2,225,806
  \item \textmark{Number of words}: 413,338
\end{itemize}

\mode<presentation>{\vfill\pause}
\begin{center}
\begin{tabular}{|l|r|r|}
\hline
\textmark{Implementation} & \textmark{Average time (ms)} & \textmark{Speedup} \\
\hline
\hline
\textgood{Classic} & 13.756 & 1.00 \\
\hline
\textgood{Classic STL} & 8.004 & 1.72 \\
\hline
\textgood{Parallel STL} & 1.824 & 7.54 \\
\hline
\textgood{Ranges} & 12.789 & 1.08 \\
\hline
\textgood{Parallel Ranges} & 1.179 & 11.67 \\
\hline
\end{tabular}
\end{center}
\end{frame}